\color{unidarkgreen}{\subsection{Логика очувствления}}
\color{black}

\begin{enumerate}[label=\arabic{section}.\arabic{subsection}.\arabic{enumi}, labelsep=4pt, leftmargin=0pt, itemindent=57pt, itemsep=0pt, parsep=5pt] % если что можно крутить itemsep


%%%%%%%%%%%%%%%%%%%%%%%%%%%%%%%%%%%%%%%%%%%%%%%%%% Логика очувствления %%%%%%%%%%%%%%%%%%%%%%%%%%%%%%%%%%%%%%%%%%%%%%%%%%
\item Общие данные

\begin{enumerate}[label=\arabic{section}.\arabic{subsection}.\arabic{enumi}.\arabic*, labelsep=4pt, leftmargin=0em, itemindent=65pt, parsep=0pt]

\item
В устройстве предусмотрено два режима очувствления ДЗО: оперативный и автоматический.
Оперативный режим вводится наличием внешнего сигнала \textit{«Оперативный ввод очувствления»} и применяется
для обеспечения требуемой чувствительности ДЗО при выводе в ремонт мощных присоединений.
Автоматический режим вводится наличием внешнего сигнала \textit{«Нормальный режим очувствления»} или от
программного ключа \textbf{«Очувст.»} и предназначен для повышения надежности отключения выключателей
присоединений при работе ДЗО, в том числе при АПВ шин и при опробовании.

\item
Режим очувствления, вводимый срабатыванием ДЗО, сохраняется на всё время отключения СШ. При
наличии сигнала \textit{«Фикс. откл СШ»} режим очувствления удерживается до момента возможного повторного
включения СШ при АПВ. Сброс очувствления до повторного включения осуществляется при наличии сигнала
\textit{«Работа АПВ»}.

\item
Переход в режим очувствления при опробовании осуществляется при отсутствии напряжения на СШ.
Очувствление в данном случае сохраняется на всё время опробования.

\item
При АПВ шин режим очувствления вводится на время включения выключателя и последующего
отключения в случае неуспешного АПВ. Для исключения излишнего очувствления при постепенном включении
всех выключателей предусмотрена возможность вывода режима очувствления при успешном АПВ по
истечению выдержки времени, задаваемой уставкой \textbf{«Точ. АПВ»}. Вывод режима осуществляется
программным ключом \textbf{«Вывод очув. при успеш. АПВ»}.

\medskip

\begin{figure}[H]
\centering
\includegraphics[width=0.9\textwidth,height=0.9\textheight,keepaspectratio]{\fbpath/661.1771 Логика очувствления/img/Log_ochuv.pdf}
\captionof{figure}{Функциональный блок логики очувствления}\label{pic_log_ochuv:Log_ochuv}
\end{figure}

\medskip

\item
Уставки функции логики очувствления представлены в таблице \ref{log_ochuv:tbl1}.

\small
\begin{longtable}{|>{\centering\arraybackslash}m{5.3cm}|>{\centering\arraybackslash}m{3.3cm}|>{\centering\arraybackslash}m{4.2cm}|>{\centering\arraybackslash}m{1.8cm}|>{\centering\arraybackslash}m{1cm}|}
% Первая страница: полный заголовок
\caption{Перечень уставок функции логики очувствления\hfill\vspace{-0.5\baselineskip}}\label{log_ochuv:tbl1}\\
\hline
\rowcolor{gray!30}
Параметр & Условное обозначение на схеме & Значение/ Диапазон & Единица измерения & Шаг \\
\hline
\endfirsthead

% Промежуточные страницы: «Продолжение таблицы»
\caption*{\hspace{3pt}\emph{Продолжение таблицы \ref{log_ochuv:tbl1}}\hfill\ }\vspace{-0.5\baselineskip} \\
\hline
\rowcolor{gray!30}
Параметр & Условное обозначение на схеме & Значение/ Диапазон & Единица измерения & Шаг \\
\hline
\endhead

%===m>ДЗО|-
%===m
\end{longtable}
\normalsize	

\end{enumerate}

\end{enumerate}
